% GENERATED FILE. Edit canonical lessons, then run npm run generate:books.
% canonical-source-hash: fnv1a64:4fabd3a9486f235b
% canonical-lessons: PA-W08-pha, PA-W08-pairin-bindi, PA-W08-hora, PA-W08-phone-label, PA-W08-digit-zero, PA-W08-phone-a, PA-W08-phone-b, PA-W08-digit-recognition, PA-W08-phone-select, PA-W08-phone-supported, PA-W08-phone-grouping, PA-W08-phone-delayed, PA-W08-phone-dictation, PA-W08-phone-repair, PA-W08-phone-no-model

\chapter{Build One Fictional Phone Field}
\label{ch:pa-form-phone-runway}
\hlchaptermodality{30}

\begin{chapteropening}
\textbf{By the end of this chapter:} \emph{I can recognise the Gurmukhi phone label and independently select and write one of two taught fictional phone values digit by digit from a nonverbal cue.}

The learner fills one fictional phone line from a remembered nonverbal cue with no romanization, Latin digits, or copyable Gurmukhi answer, then checks selection, six digit positions, grouping, and placement separately.
\end{chapteropening}

\section[pha]{\pa{ਫ} — one new base shape}

\label{lesson:PA-W08-pha}



\begin{quote}
Write the familiar \textbf{\pa{ਪ}} once. The new base begins similarly but is not the same letter.
\end{quote}

\subsection*{Script — one letter}
Trace \textbf{\pa{ਫ}} slowly twice. This is \emph{pha}. Stop before adding any dot or vowel mark.

\subsection*{Your turn}
Cover the model, write \textbf{\pa{ਫ}} once, then compare only its shape.

\begin{tcolorbox}[breakable,colback=teal!4,colframe=teal!35!black,title={Before you move on}]
One new base is enough.
\end{tcolorbox}
\section[fa]{\pa{ਫ਼} — place one dot below}

\label{lesson:PA-W08-pairin-bindi}



\begin{quote}
Write \textbf{\pa{ਫ}} once and leave room directly below it.
\end{quote}

\subsection*{Script — one mark}
Add the small dot below: \textbf{\pa{ਫ} + \pa{਼} = \pa{ਫ਼}}. This below-letter dot is \emph{pairin bindi}; here it changes \emph{pha} to the borrowed sound \emph{fa}. It is not the above-letter bindi already used in \textbf{\pa{ਹਾਂ}}.

\subsection*{Your turn}
Copy \textbf{\pa{ਫ਼}} once. Point to the dot below before you stop.

\begin{tcolorbox}[breakable,colback=teal!4,colframe=teal!35!black,title={Before you move on}]
Only the base and its below-letter dot are scored here.
\end{tcolorbox}
\section[hora]{\pa{ੋ} — add the long-o mark}

\label{lesson:PA-W08-hora}



\begin{quote}
Write \textbf{\pa{ਫ਼}} once.
\end{quote}

\subsection*{Script — one vowel mark}
The vowel mark \textbf{\pa{ੋ}}, called \emph{hora}, sits above and to the right of its base. Add it to the known base: \textbf{\pa{ਫ਼} + \pa{ੋ} = \pa{ਫ਼ੋ}}.

\subsection*{Your turn}
Copy \textbf{\pa{ਫ਼ੋ}} once. Check that the dot stays below and hora stays above-right.

\begin{tcolorbox}[breakable,colback=teal!4,colframe=teal!35!black,title={Before you move on}]
Stop after this one vowel mark.
\end{tcolorbox}
\section[fon]{\pa{ਫ਼ੋਨ} — assemble the phone label}

\label{lesson:PA-W08-phone-label}



\begin{quote}
Touch the pieces in order: \textbf{\pa{ਫ}}, the dot below, \textbf{\pa{ੋ}}, \textbf{\pa{ਨ}}.
\end{quote}

\subsection*{Script — one label}
Join the pieces as \textbf{\pa{ਫ਼ੋਨ}}. On this fictional beginner form, \textbf{\pa{ਫ਼ੋਨ}} labels the phone-number line. Copy the label once and leave its value blank.

\subsection*{Your turn}
Point to \textbf{\pa{ਫ਼ੋਨ}} and say what information belongs there. Do not add digits.

\begin{tcolorbox}[breakable,colback=teal!4,colframe=teal!35!black,title={Before you move on}]
No real contact information is requested.
\end{tcolorbox}
\section[zero]{\pa{੦} — add the digit zero}

\label{lesson:PA-W08-digit-zero}



\begin{quote}
Write the three known Gurmukhi digits: \textbf{\pa{੧}, \pa{੨}, \pa{੫}}.
\end{quote}

\subsection*{Script — one digit}
Trace \textbf{\pa{੦}} twice. It is the Gurmukhi digit zero. Phone values are identifiers, so a leading zero stays visible.

\subsection*{Your turn}
Cover the model and write \textbf{\pa{੦}} once.

\begin{tcolorbox}[breakable,colback=teal!4,colframe=teal!35!black,title={Before you move on}]
All four digit shapes needed by this chapter are now known.
\end{tcolorbox}
\section[zero-one-two-two-five-one]{\pa{੦੧੨੨੫੧} — build fictional value A}

\label{lesson:PA-W08-phone-a}



\begin{quote}
Point to the digits as they are named: zero, one, two, five.
\end{quote}

\subsection*{Script — six places}
Fictional practice value A is \textbf{\pa{੦੧੨੨੫੧}}. Touch and say each place from left to right: \textbf{\pa{੦} · \pa{੧} · \pa{੨} · \pa{੨} · \pa{੫} · \pa{੧}}. This short value is only curriculum data; it is not a real or callable number.

\subsection*{Your turn}
Copy A once, pausing after the third digit. Do not add grouping yet.

\begin{tcolorbox}[breakable,colback=teal!4,colframe=teal!35!black,title={Before you move on}]
Only one six-place value is active.
\end{tcolorbox}
\section[zero-two-five-one-two-five]{\pa{੦੨੫੧੨੫} — build fictional value B}

\label{lesson:PA-W08-phone-b}



\begin{quote}
Touch value A, \textbf{\pa{੦੧੨੨੫੧}}, from left to right once.
\end{quote}

\subsection*{Script — a second six-place value}
Fictional practice value B is \textbf{\pa{੦੨੫੧੨੫}}. Touch each place: \textbf{\pa{੦} · \pa{੨} · \pa{੫} · \pa{੧} · \pa{੨} · \pa{੫}}. It uses only the four introduced digit shapes.

\subsection*{Your turn}
Copy B once, pausing after the third digit. Do not add grouping yet.

\begin{tcolorbox}[breakable,colback=teal!4,colframe=teal!35!black,title={Before you move on}]
The bounded bank now contains exactly A and B.
\end{tcolorbox}
\section[zero, one, two, five]{Read one place at a time}

\label{lesson:PA-W08-digit-recognition}



\begin{quote}
Shuffle four cards marked \textbf{\pa{੦}, \pa{੧}, \pa{੨}, \pa{੫}}. Name each card before writing.
\end{quote}

\subsection*{You'll want to know — locate, do not guess}
In A, circle the second digit. In B, circle the fifth digit. Read the complete values only after locating those positions.

\subsection*{Your turn}
Hear: \textbf{zero, one, two, two, five, one}. Write one digit after each word, then compare with A.

\begin{tcolorbox}[breakable,colback=teal!4,colframe=teal!35!black,title={Before you move on}]
No unintroduced numeral is scored.
\end{tcolorbox}
\section[two target-script selectors]{Choose A or B before writing}

\label{lesson:PA-W08-phone-select}



\begin{quote}
Read the closed bank once: A \textbf{\pa{੦੧੨੨੫੧}}; B \textbf{\pa{੦੨੫੧੨੫}}.
\end{quote}

\subsection*{You'll want to know — two fixed selectors}
For this fictional practice only, \textbf{\pa{ਕ}} means A and \textbf{\pa{ਖ}} means B. Point to the matching value before your pen moves. The selector is not part of the phone field.

\subsection*{Your turn}
For \textbf{\pa{ਖ}}, select B. For \textbf{\pa{ਕ}}, select A. Selection and digit writing are separate checks.

\begin{tcolorbox}[breakable,colback=teal!4,colframe=teal!35!black,title={Before you move on}]
The selectors carry no special meaning outside this exercise.
\end{tcolorbox}
\section[phone field]{One supported phone field}

\label{lesson:PA-W08-phone-supported}



\begin{quote}
Point to \textbf{\pa{ਫ਼ੋਨ}}. Read the visible bank: A \textbf{\pa{੦੧੨੨੫੧}}; B \textbf{\pa{੦੨੫੧੨੫}}.
\end{quote}

\subsection*{Writing — bank visible}
The selector is \textbf{\pa{ਕ}}. Complete \textbf{\pa{ਫ਼ੋਨ}: \_\_\_\_\_\_\_\_\_\_} with A while the bank remains visible. Write one digit per place.

\subsection*{Your turn}
Point under each copied digit and compare all six places left to right.

\begin{tcolorbox}[breakable,colback=teal!4,colframe=teal!35!black,title={Before you move on}]
This supported copy does not earn independent-writing credit.
\end{tcolorbox}
\section[three digits, space, three digits]{Three digits, one space, three digits}

\label{lesson:PA-W08-phone-grouping}



\begin{quote}
Touch A in six places: \textbf{\pa{੦} · \pa{੧} · \pa{੨} · \pa{੨} · \pa{੫} · \pa{੧}}.
\end{quote}

\subsection*{Script — one grouping boundary}
Insert one clear space after the third digit: \textbf{\pa{੦੧੨} \pa{੨੫੧}}. The space helps the eye track the practice value; it does not remove or reorder a digit. B groups as \textbf{\pa{੦੨੫} \pa{੧੨੫}}.

\subsection*{Your turn}
Copy grouped B once. Count three digits, make one space, then count three more.

\begin{tcolorbox}[breakable,colback=teal!4,colframe=teal!35!black,title={Before you move on}]
This book keeps one grouping inside the bounded practice bank; it is not a claim about every real-world phone format.
\end{tcolorbox}
\section[delayed phone-field entry]{Hide the bank, then write B}

\label{lesson:PA-W08-phone-delayed}



\begin{quote}
Study \textbf{\pa{ਫ਼ੋਨ}: \pa{੦੨੫} \pa{੧੨੫}} for five seconds. Say each digit while pointing.
\end{quote}

\subsection*{Writing — short delay}
Hide the bank. Wait five seconds. For \textbf{\pa{ਖ}}, complete one blank \textbf{\pa{ਫ਼ੋਨ}} field from memory.

\subsection*{Your turn}
Uncover the model only after writing. Compare digit selection, order, and the single group space separately.

\begin{tcolorbox}[breakable,colback=teal!4,colframe=teal!35!black,title={Before you move on}]
The model is hidden during the attempt.
\end{tcolorbox}
\section[heard digit string]{Hear six digits and write them}

\label{lesson:PA-W08-phone-dictation}



\begin{quote}
Without a number model, write the four introduced digits when you hear: zero, one, two, five.
\end{quote}

\subsection*{Writing — digit by digit}
Listen twice: \textbf{zero — one — two | two — five — one}. The pause marks the one group space. Write Gurmukhi digits; do not write English numerals or words.

\subsection*{Your turn}
Read your six digits back left to right. Check the space only after checking all six positions.

\begin{tcolorbox}[breakable,colback=teal!4,colframe=teal!35!black,title={Before you move on}]
Romanization and Latin digits do not earn the Gurmukhi writing score.
\end{tcolorbox}
\section[four repair passes]{Repair only one dimension}

\label{lesson:PA-W08-phone-repair}



\begin{quote}
Name four separate checks: cue selection, six digit positions, middle group space, and placement on the \textbf{\pa{ਫ਼ੋਨ}} line.
\end{quote}

\subsection*{Your turn — bounded repair}
For \textbf{\pa{ਕ}}, inspect \textbf{\pa{ਫ਼ੋਨ}: \pa{੦੧੨੨} \pa{੫੧}}. The selected value and digit order are right; only the group boundary is wrong. Repair it as \textbf{\pa{੦੧੨} \pa{੨੫੧}} without rewriting correct digits.

\subsection*{Your turn}
Run the checks in order. Stop after the first differing dimension is repaired.

\begin{tcolorbox}[breakable,colback=teal!4,colframe=teal!35!black,title={Before you move on}]
Do not replace the whole answer when one dimension differs.
\end{tcolorbox}
\section[phone field]{One no-model phone field}

\label{lesson:PA-W08-phone-no-model}



\begin{quote}
Close the previous lessons. There is no value bank, support-language label, Latin-digit version, or copyable Gurmukhi answer below.
\end{quote}

\subsection*{Writing — independent controlled choice}
Complete the requested fictional field:

\begin{quote}
\pa{ਖ} — \textbf{\pa{ਫ਼ੋਨ}: \_\_\_\_\_\_\_\_\_\_}
\end{quote}

Choose the known value attached to the circle and write it from memory in Gurmukhi digits. Do not reopen an earlier lesson until the attempt is complete.

\subsection*{Your turn — separate repair passes}
After writing, compare cue selection, six digit positions, grouping, and field placement. Repair only the first differing dimension.

\begin{tcolorbox}[breakable,colback=teal!4,colframe=teal!35!black,title={Before you move on}]
This closes only the untimed phone-field ladder. Date, integration, and A1 readiness remain separate dependency-linked work.
\end{tcolorbox}
